\subsection{Marco teórico}

El desarrollo de un middleware de orquestación para realizar anotaciones automáticas en ELAN requiere integrar fundamentos de lingüística computacional, procesamiento de videos, arquitecturas de software y ejecución controlada de modelos de Inteligencia Artificial. En este capítulo se presentan los conceptos necesarios para comprender el problema técnico y proporcionar el contexto teórico indispensable para entender el origen y la justificación de la solución desarrollada.
El marco teórico se organiza en siete bloques fundamentales, siendo el primero la revisión de la anotación multimodal y el papel de ELAN en la construcción de corpus audiovisuales. Luego, se analiza el uso de reconocimiento automático para apoyar tareas de segmentación y etiquetado de videos. A continuación, se presentan las arquitecturas de integración entre aplicaciones de escritorio heredadas y servicios externos, dando énfasis en el uso del patrón Sidecar. Posteriormente, se estudian los servicios REST y los contratos JSON como mecanismos de comunicación entre sistemas. También se revisa el uso de contenedores como la base conceptual para aislar dependencias y ejecutar modelos en entornos reproducibles. Finalmente, se aborda el tema de la orquestación de modelos de inteligencia artificial y las consideraciones de interoperabilidad, trazabilidad y mantenibilidad que justifican la solución propuesta.\vspace{1\baselineskip}

\subsubsection{Herramientas de anotación multimodal}

Para comenzar, la anotación multimodal consiste en asociar descripciones lingüísticas, gestuales, acústicas o visuales con intervalos específicos de tiempo de una señal de audio o video. En investigaciones realizadas sobre interacción humana, gestos, conversación y lenguas de señas, la anotación no se la encasilla simplemente como transcribir palabras, sino que se interpreta como registrar eventos sincronizados en el tiempo, relaciones entre capas de análisis y metadatos que permiten recuperar, comparar y reutilizar un corpus. Por esta razón, las herramientas de anotación deben ofrecer precisión temporal, soporte para múltiples niveles de descripción y capacidad de vincular anotaciones con diferentes tipos de archivos multimedia \cite{wittenburg2006elan}.

ELAN utiliza archivos nativos en formato EAF para representar anotaciones mediante estructuras temporales que vinculan valores textuales con marcas de inicio y fin. Hablando conceptualmente, los archivos EAF permiten que una anotación sea tratada como un segmento de señal multimedia, lo cual es indipensable cuando se analizan secuencias audiovisuales. La principal utilidad de este formato no se encuentra únicamente en almacenar texto, sino en mantener la correspondencia entre el evento observado y su posición temporal. En tareas de anotación asistida, esta correspondencia se convierte en el punto de enlace entre la salida de un modelo que realice segmentación automática y la revisión posterior de un especialista humano.

ELAN se destaca en investigación multimodal por su flexibilidad y énfasis en la exactitud temporal \cite{wittenburg2006elan}. Además, se han propuesto extensiones para integrar detectores de sonido e imagen dentro del flujo de anotación, como se evidenció en el proyecto AVATecH, el cual está orientado a incorporar componentes de reconocimiento de patrones en procesos de anotación semiautomática \cite{auer2010elan}. Esto es relevante para el presente TIC porque muestra que la anotación asistida no reemplaza el juicio del investigador, sino que introduce mecanismos computacionales para proponer segmentos o etiquetas que deben ser revisados.

En el contexto de la Lengua de Señas Ecuatoriana (LSEC), la anotación audiovisual adquiere una importancia particular. La LSEC constituye una lengua visual-gestual usada por la comunidad sorda ecuatoriana y su documentación demanda herramientas que permitan conservar la relación entre movimiento, temporalidad y significado. Estudios han señalado la relevancia de la LSEC para procesos de inclusión educativa y comunicación accesible \cite{ureta2022lsec_inclusion}. Asimismo, trabajos previos de la Escuela Politécnica Nacional han explorado el uso de técnicas de visión artificial y aprendizaje profundo para interpretar frases de LSEC, lo que evidencia un interés creciente por relacionar corpus visuales con modelos computacionales \cite{guevara2022lsec_interprete}.

La construcción de un corpus audiovisuales para LSEC exige una etapa de anotación cuidadosa. Una glosa no representa necesariamente toda la riqueza lingüística de una seña, pero funciona como una etiqueta de trabajo que permite indexar, buscar y entrenar modelos sobre segmentos temporales. En consecuencia, la anotación manual sigue siendo una referencia esencial para validar resultados automáticos. La asistencia computacional puede reducir carga repetitiva o proponer intervalos iniciales, pero la decisión final sobre la corrección lingüística de una etiqueta depende del criterio humano y de la calidad del corpus disponible.

Por último, los sistemas de anotación semiautomática deben considerar que el material audiovisual puede ser heterogéneo. La variación en iluminación, encuadre, velocidad de ejecución, morfología del señante, cámara, fondo y disponibilidad de metadatos incide en la confiabilidad de los resultados. Por esta razón, el marco teórico para esta solución debe contemplar tanto la estructura de los archivos de anotación como las limitaciones de los modelos automáticos. El objetivo conceptual no es convertir a ELAN en un sistema de inteligencia artificial, sino permitir que una herramienta de anotación consolidada pueda comunicarse con servicios externos especializados.

\subsubsection{Reconocimiento automático en procesos de anotación}
------------------------------------------------
El reconocimiento automático de señas o segmentos en videos consiste en extraer patrones de una señal visual para producir inferencias sobre eventos, acciones, gestos o categorías. En un proceso de anotación lingüística, estas inferencias suelen expresarse como segmentos temporales, etiquetas candidatas, puntuaciones de confianza o listas de alternativas. El reconocimiento automático se diferencia de la anotación manual debido a que opera usando modelos estadísticos o de aprendizaje profundo, los cuales están entrenados sobre datos previos, mientras que la anotación manual depende de la observación y el criterio de un experto en el área.

En archivos multimedia, la segmentación temporal busca identificar intervalos donde ocurren eventos relevantes. En lenguaje de señas, esto es equivalente a detectar cuándo inicia y finaliza un gesto, una seña o una frase . En cambio, la clasificación intenta asignar una etiqueta al segmento identificado. Estas dos tareas pueden ejecutarse en un único modelo o se puede hacer uso de una cadena de modelos, si se implementauna cadena es importante recalcar que primero se detectan intervalos y luego se clasifican. Separar ambas funciones permite tratar de forma independiente el problema temporal y el problema linguístico. Esta separación resulta conveniente cuando el corpus es limitado o cuando los errores de segmentación y clasificación deben analizarse por separado.

La anotación asistida se sitúa entre la automatización completa y el trabajo manual. El sistema propone segmentos o etiquetas, y el experto se encarga de validar, corregir o descarta los resultados obtenidos. Este enfoque es importante en la creación de corpus audiovisuales porque los modelos pueden acelerar tareas repetitivas, pero no eliminan la necesidad de que un experto revise los resultados. La validación humana sigue siendo necesaria por la variabilidad natural de las señas, la posible ambigüedad entre las glosas, la dependencia del contexto y la presencia de rasgos no manuales que pueden no ser capturados por un modelo particular.

Desde el punto de vista de representación, un modelo puede trabajar directamente con frames de video o con características principales extraídas del video. El procesamiento basado en frames conserva información visual amplia, pero puede aumentar el costo computacional y requerir arquitecturas capaces de manejar dimensiones espaciales y temporales. El procesamiento basado en puntos clave convierte el video a coordenadas corporales o manuales, lo cual puede ayudar a disminuir ruido de fondo y facilitar el análisis secuencial. La herramienta MediaPipe se ubica en esta segunda línea dado que proporciona una infraestructura para construir pipelines de percepción y extraer landmarks corporales, faciales o de manos en aplicaciones de visión por computadora \cite{lugaresi2019mediapipe}.

Los modelos secuenciales son relevantes cuando la salida depende del orden temporal de la señal. Las redes LSTM fueron propuestas para modelar dependencias en secuencias y mitigar problemas asociados al aprendizaje de relaciones de largo plazo \cite{hochreiter1997lstm}. En procesamiento de señas, esta capacidad es útil porque la identidad de una seña no depende solo de una postura estática, sino del movimiento previamente realizado. Las variantes bidireccionales en redes LSTM permiten considerar información contextual antes y después de una posición temporal, lo que resulta útil para segmentar eventos dentro de una secuencia completa dado que nos permite procesar los datos de una manera más efectiva.

El esquema BIO, utilizado originalmente en tareas de segmentación de texto, codifica posiciones mediante etiquetas que indican inicio, interior o exterior de una unidad \cite{ramshaw1995text_chunking}. En un video, el mismo principio puede adaptarse para marcar frames como inicio de segmento, continuación de segmento o fondo. Esta formulación permite convertir una predicción por frame en intervalos temporales. Sin embargo, requiere postprocesamiento para corregir secuencias inconsistentes, unir fragmentos cercanos, eliminar segmentos demasiado cortos y calcular valores de confianza.

Los modelos basados en Transformer introdujeron mecanismos de atención que permiten ponderar relaciones entre elementos de una secuencia sin depender exclusivamente de recurrencia \cite{vaswani2017attention}. En clasificación de glosas, una arquitectura Transformer puede procesar una secuencia de puntos clave y producir una distribución de probabilidad sobre un vocabulario. Esto no implica que el modelo comprenda todos los aspectos lingüísticos de la seña, implica que aprende patrones asociados a etiquetas observadas durante el entrenamiento. Por ello, el tamaño y la diversidad del corpus, la calidad de las anotaciones y la cobertura del vocabulario influyen directamente en la utilidad de este clasificador.

La disponibilidad de datos útiles es una limitante crítica en lenguaje de señas. Cuando un corpus contiene pocas muestras por glosa, segmentaciones no balanceadas o anotaciones incompletas. El modelo puede aprender patrones específicos del conjunto disponible y generalizar de forma limitada. Por esta razón, las salidas automáticas del modelos deben presentarse como apoyo y no como un sustituto de la anotación realizada por un experto. En términos de diseño, esto justifica que una respuesta de reconocimiento incluya no solo una etiqueta final, sino también confianza, alternativas y trazabilidad suficiente para que un experto evalúe el resultado.

Este proyecto sustenta conceptualmente esta separación entre segmentación y clasificación mediante un pipeline que realiza extracción de puntos clave, segmentación BIO, clasificación de glosas y generación de segmentos con campos temporales. En esta sección, cada elemento mencionado se presentan como fundamentos de reconocimiento automático.

\subsubsection{Arquitecturas de integración}

Las aplicaciones de escritorio consolidadas suelen concentrar flujos de trabajo especializados, interfaces maduras y formatos de archivo propios. Sin embargo, su arquitectura puede no estar preparada para incorporar dependencias modernas de inteligencia artificial, especialmente cuando estas dependen de lenguajes, bibliotecas nativas, aceleración por hardware o entornos de ejecución distintos. Bajo este contexto, integrar modelos directamente dentro de la aplicación principal puede aumentar el acoplamiento, dificultar el mantenimiento y comprometer la estabilidad del software.

Una alternativa a nivel de arquitectura de software consiste en implementar una capa intermedia (Middleware). En términos generales, el middleware actuará como un intermediario entre sistemas que no comparten la misma interfaz, lenguaje, protocolo o ciclo de vida. Su función no se limita a reenviar mensajes, sino también a validar datos, traducir contratos, gestionar errores, seleccionar servicios, registrar eventos y aislar dependencias. En el presente proyecto, este concepto se relaciona con la necesidad de conectar ELAN, construido sobre el ecosistema Java, con servicios de inteligencia artificial implementados en Python y bibliotecas como PyTorch.

El patrón empleado, en este caso sidecar describe una forma de organizar componentes auxiliares alrededor de una aplicación principal. En arquitecturas basadas en contenedores, este patrón permite que un servicio acompañante proporcione funcionalidades de soporte sin modificar de forma invasiva la lógica central de la aplicación \cite{burns2016container_patterns}. En el contexto de una aplicación de escritorio, el principio se mantiene porque el componente auxiliar ejecuta tareas especializadas y expone una interfaz controlada, mientras que la aplicación principal conserva su flujo de trabajo y únicamente delega las operaciones externas necesarias.

La aplicación principal, en este caso ELAN, mantiene el rol de herramienta de anotación y visualización. El middleware implementado se concibe como una capa de orquestación local encargada de recibir solicitudes, validar parámetros, seleccionar modelos y devolver resultados estructurados. La separación implementada permite que los modelos de inteligencia artificial evolucionen sin exigir que ELAN incorpore de forma nativa todas sus dependencias. Además, reduce significativamente el riesgo de que una falla de inferencia o una incompatibilidad de bibliotecas afecte directamente al proceso de anotación.

La integración entre aplicaciones de escritorio y servicios externos puede realizarse mediante archivos intermedios, llamadas de sistema, sockets, servicios HTTP u otros mecanismos de comunicación. La elección depende de requisitos como portabilidad, facilidad de depuración, seguridad, rendimiento y soporte de bibliotecas. Para un entorno local, un servicio HTTP sobre la IP local ofrece una frontera clara entre diferentes procesos y permite reutilizar herramientas comunes de prueba, documentación y validación. Esta decisión también favorece la independencia entre lenguajes, dado que la aplicación Java y el servicio Python solo necesitan seguir el contrato de comunicación establecido.

El desacoplamiento es un principio central en este tipo de arquitectura. Cuando el cliente, el middleware y los modelos comparten contratos explícitos, cada parte puede evolucionar sin impactar sobre las demás de forma directa. Para ELAN es indistinto conocer si la inferencia se ejecuta mediante un runner simulado, un runner nativo de PyTorch o un runner desplegado en contenedores. Del mismo modo el middleware puede incorporar nuevos modelos siempre que estos respeten los contratos de entrada y salida previamente definidos para el sistema.

La arquitectura de integración también debe considerar la experiencia del usuario final dado que, en un flujo de anotación, el investigador espera seleccionar material audiovisual, activar una tarea de reconocimiento y recibir segmentos editables dentro del entorno de trabajo. Por esta razón, el middleware no impone un cambio conceptual sobre la tarea lingüística, su valor consiste en ocultar complejidad técnica y devolver resultados compatibles con las estructuras de anotación nativa de ELAN. De esta manera, la innovación se concentra en la interoperabilidad y no en reemplazar o modificar la herramienta original.

\subsubsection{Servicios REST y contratos de comunicación}

REST, propuesto por Fielding como estilo arquitectónico para sistemas distribuidos, define principios para organizar recursos, representaciones y operaciones sobre una interfaz uniforme \cite{fielding2000rest}. En un servicio REST, las funcionalidades se exponen mediante endpoints, pueden ser utilizadas con métodos HTTP, y las respuestas se comunican con códigos de estado y representaciones estructuradas. Aunque REST no establece un formato de datos único, JSON se utiliza ampliamente por su compatibilidad con múltiples lenguajes y por su facilidad para representar objetos, listas, cadenas, números y valores booleanos.

En una integración entre ELAN y un middleware de inteligencia artificial, la API REST permite convertir una solicitud de anotación en un mensaje detallado de forma explícita. El cliente puede enviar la ruta del video, la configuración del modelo, el tier destino, parámetros de segmentación y preferencias de ejecución. El middleware responde con estado del job, información del medio, segmentos temporales y trazas de ejecución. Esta forma de comunicación evita acoplar directamente objetos internos de Java con objetos internos de Python, también permite probar el servicio mediante clientes HTTP independientes antes de integrarlo con la aplicación de escritorio final.

Los contratos de comunicación son esenciales para garantizar interoperabilidad. Un contrato define campos obligatorios, tipos de datos, rangos permitidos, estructura de respuesta y errores esperados. En el sistema analizado, los campos \texttt{start\_ms}, \texttt{end\_ms}, \texttt{label} y \texttt{confidence} representan el núcleo de un segmento anotable. Los dos primeros delimitan el intervalo temporal, el tercero contiene la etiqueta propuesta y el cuarto expresa una medida de confianza. Esta estructura es suficiente para que un cliente pueda convertir resultados automáticos en anotaciones temporales visibles en la linea de tiempo en la que el linguista trabaja.

La validación de esquemas reduce errores tempranos. FastAPI proporciona una infraestructura para declarar endpoints, asociar modelos de datos y generar documentación OpenAPI a partir de tipos de Python \cite{fastapi2026features}. Pydantic, por su parte, permite definir modelos de validación y conversión de datos, de modo que un payload incorrecto pueda rechazarse antes de ejecutar lógica de negocio \cite{pydantic2026models}. En un middleware de inferencia, esto es importante dado que reduce el poder computacional innecesario al evitar la ejecución de modelos cuando faltan campos, los tipos no coinciden o los parámetros no cumplen restricciones mínimas .

Los códigos de estado y las respuestas de error forman parte del contrato. Una API robusta no debe limitarse a devolver datos cuando todo funciona, también debe informar si el video no existe, si el modelo no está registrado, si la arquitectura no coincide, si el procesamiento CUDA fue solicitada pero no está disponible o si el vocabulario de glosas no cumple el formato esperado. Las respuestas estructuradas facilitan la depuración y permiten que el cliente presente mensajes controlados sin cerrar la aplicación principal. El correcto manejo de códigos de estado y respuestas de error resultan especialmente útiles cuando el servicio se integra con un flujo de anotación interactivo dado que nos permiten depurar, realizar pruebas y mapear errores en la implementación.

AVATecH aporta un antecedente relevante para la comunicación entre ELAN y reconocedores externos. Su especificación técnica describe mecanismos para que extensiones de reconocimiento intercambien parámetros, archivos y resultados con ELAN \cite{auer2014avatech}. Aunque la tecnología y los modelos actuales han evolucionado, el principio se mantiene dado que la herramienta de anotación requiere una frontera bien definida para invocar procesos externos y recibir anotaciones candidatas. La traducción conceptual de un protocolo heredado hacia una API REST moderna debe preservar la semántica temporal y la estructura de los resultados para no generar problemas de compatibilidad de formatos.

Los metadatos también influyen en la interoperabilidad. CMDI fue propuesto como un marco basado en componentes y categorías de datos para mejorar la consistencia de metadatos en recursos lingüísticos \cite{broeder2010cmdi}. En un sistema de anotación asistida, los metadatos permiten describir modelos, entradas, salidas, parámetros y recursos asociados. Esto facilita que los resultados generados por un componente automático no queden aislados, sino que puedan relacionarse con el corpus, la herramienta de anotación y la configuración de ejecución.

\subsubsection{Contenedores y ejecución controlada de modelos}

La contenedorización permite empaquetar una aplicación junto con dependencias, bibliotecas y configuración de ejecución en una unidad portable. Docker describe los contenedores como entornos aislados que ejecutan procesos sobre el sistema anfitrión y que pueden construirse a partir de imágenes reproducibles en cualquier sistema \cite{docker2026overview}. Para sistemas de inteligencia artificial, este enfoque resulta atractivo porque los modelos suelen depender de versiones específicas de Python, bibliotecas numéricas, frameworks de aprendizaje profundo, controladores y recursos de hardware.

En una arquitectura que integra ELAN con modelos de inteligencia artificial, el hacer uso de contenedores ofrece una separación clara entre la herramienta de anotación y el entorno de inferencia. ELAN no necesita instalar directamente PyTorch, OpenCV, MediaPipe o dependencias nativas. El modelo puede ejecutarse en un entorno controlado y exponer resultados mediante el middleware. Esta separación reduce conflictos de versiones y permite reproducir el comportamiento del servicio en distintos equipos, siempre que se mantenga la misma imagen o configuración de ejecución.

La ejecución nativa y la ejecución contenerizada representan dos estrategias distintas. En la ejecución nativa, el servicio de middleware carga las bibliotecas directamente desde el entorno Python local. Esta modalidad puede ser ampliamente utilizada en el desarrollo dado que permite depurar con menor sobrecarga. En la ejecución contenerizada, el modelo se ejecuta en un entorno aislado, lo que mejora la portabilidad y el control de dependencias, pero introduce cargas considerables al momento de realizar el montaje de archivos, comunicación con el contenedor, acceso a GPU, límites de recursos y ciclo de vida del proceso.

El uso de contenedores no elimina la necesidad de contratos. Un modelo contenerizado debe recibir entradas en un formato acordado y devolver salidas compatibles con el middleware. Por elloel hacer uso de contenedores debe entenderse como una estrategia de aislamiento y reproducibilidad, no como un sustituto de una arquitectura de comunicación. El middleware sigue siendo responsable de validar solicitudes, seleccionar el runner adecuado, controlar errores y convertir resultados a una representación que ELAN pueda utilizar.

La ejecución controlada de modelos también se relaciona con recursos de cómputo. Los modelos de video pueden requerir memoria considerable, aceleración por GPU o tiempos de inferencia considerablemente altos. Un sistema de orquestación debe contemplar parámetros como dispositivo preferido, tiempo máximo de ejecución y manejo de fallos. No obstante, las métricas de consumo, latencia y rendimiento también son fundamentales para evaluar el rendimiento del sistema implemetado.

Desde el punto de vista académico y técnico, el hacer uso de contenedores mejora la reproducibilidad en diferentes enotorno y separar responsabilidades. La herramienta de anotación conserva su entorno nativo, el middleware coordina solicitudes y los modelos se ejecutan bajo condiciones definidas. Esta organización es coherente con el principio de mantener bajo acoplamiento entre componentes y de introducir cambios técnicos sin alterar innecesariamente la experiencia del usuario final.

\subsubsection{Orquestación de modelos de inteligencia artificial}

La orquestación de modelos de inteligencia artificial consiste en coordinar las etapas necesarias para ejecutar una inferencia de forma controlada. No se limita a llamar a un modelo dado que incluye registrar modelos disponibles, validar sus artefactos, preparar entradas, seleccionar un runner, controlar dispositivo, ejecutar inferencia, transformar salidas, postprocesar resultados y registrar trazas. En un contexto de anotación, la orquestación también debe asegurar que la salida final sea interpretable como segmentos temporales o etiquetas candidatas.

Un registro de modelos es un componente que almacena información sobre modelos instalados, versiones, estado, tipo de tarea, artefactos y contratos de entrada y salida. Su utilidad radica en separar la identidad del modelo de su ejecución concreta. El cliente puede solicitar un \texttt{model\_id} y una versión, mientras que el middleware resuelve dónde se encuentran los pesos, qué runner debe utilizarse y qué restricciones aplican. Este patrón favorece la extensibilidad porque nuevos modelos pueden incorporarse sin cambiar la forma general de la solicitud.

La selección de runner es otra función central. Un runner encapsula la forma específica de ejecutar un modelo o una simulación. Puede existir un runner dummy para validar contratos, un runner nativo para cargar modelos PyTorch en el proceso local y un runner especializado para pipelines compuestos. Esta abstracción evita que el servicio de jobs tenga que conocer los detalles de cada arquitectura. En términos de diseño, el runner traduce una entrada interna normalizada a la forma que necesita el modelo y devuelve una salida interna que el postprocesador puede transformar.

PyTorch es un framework ampliamente utilizado para aprendizaje profundo, diseñado con un estilo imperativo y soporte para cómputo tensorial acelerado \cite{paszke2019pytorch}. En el presente trabajo de integración, su importancia conceptual está ligada con la ejecución de modelos de video y secuencia dentro de un entorno Python. No obstante, la existencia de PyTorch como dependencia no implica que cualquier arquitectura sea compatible automáticamente. Se requieren pesos, estructuras de modelo y formas de tensores coherentes con los contratos definidos por el middleware.

El preprocesamiento convierte el video original en una representación adecuada para el modelo. Puede incluir lectura de metadatos, muestreo de frames, conversión de espacios de color, redimensionamiento, normalización y construcción de ventanas temporales. Estas etapas son necesarias porque los modelos de aprendizaje profundo suelen requerir tensores de forma fija y rangos numéricos específicos. En el caso de modelos basados en puntos clave, el preprocesamiento también puede incluir extracción de landmarks, normalización espacial y cálculo de características dinámicas.

La inferencia produce una salida que no siempre coincide con el formato que necesita el usuario. Un segmentador puede devolver probabilidades por frame y un clasificador puede devolver logits o probabilidades por clase, por otra parte un pipeline compuesto puede devolver listas top-k. Por esta razón, el postprocesamiento es indispensable. Cuando usamos segmentación temporal, el postprocesamiento convierte probabilidades o etiquetas en intervalos con inicio, fin y confianza. Para clasificación, selecciona la etiqueta principal y puede conservar alternativas para revisión.

El pipeline BIO más clasificador de glosas responde a una división conceptual entre detección temporal y asignación léxica. La etapa BIO identifica regiones candidatas en la secuencia de keypoints, mientras que la etapa Transformer clasifica cada región contra un vocabulario de glosas. Esta organización permite que el middleware produzca segmentos cuyo campo \texttt{label} contenga una glosa propuesta y cuyo campo \texttt{predictions} incluya alternativas ordenadas. La presencia de alternativas no debe interpretarse como certeza lingüística, sino como información adicional para la revisión humana.

La trazabilidad de ejecución permite reconstruir lo que ocurre durante una solicitud. En un sistema de orquestación, una traza puede incluir runner seleccionado, dispositivo usado, modelo, versión, número de ventanas, estados atravesados y tiempos por etapa. Estas marcas no sustituyen una evaluación experimental, pero facilitan depuración, auditoría técnica y comparación controlada entre ejecuciones. Cuando un modelo falla, la traza y el código de error ayudan a distinguir si el problema está en la carga de modelo, lectura de video, extracción de keypoints, inferencia o postprocesamiento.

El diseño de la orquestación debe mantener una frontera clara entre salida interna y salida externa. Internamente, un runner puede trabajar con probabilidades por frame, tensores, logits o listas de predicciones. Externamente, el cliente requiere una respuesta estable de anotación. Esta diferencia evita que ELAN dependa de detalles específicos de PyTorch, MediaPipe o una arquitectura particular. El middleware opera como traductor entre el mundo de modelos y el mundo de anotaciones temporales.

\subsubsection{Interoperabilidad, trazabilidad y mantenibilidad}
La interoperabilidad se refiere a la capacidad de sistemas distintos para intercambiar información y utilizarla de forma significativa. En este trabajo de integración, la interoperabilidad no consiste únicamente en enviar datos desde ELAN hacia un servicio HTTP, también implica conservar el significado temporal de los segmentos, mantener etiquetas compatibles con tiers de anotación en ELAN y devolver errores comprensibles para el flujo de trabajo. La interoperabilidad efectiva requiere contratos claros, documentación técnica y una separación de responsabilidades entre cliente, middleware y modelos.

La compatibilidad con ELAN exige respetar la naturaleza temporal de la anotación. Un resultado automático solo es útil si puede ubicarse en la línea de tiempo y revisarse junto con el video. Por ello, los segmentos generados por el middleware deben expresar intervalos en milisegundos y etiquetas textuales. La generación directa de archivos EAF desde el middleware no se encuentra sustentada por el código revisado; el enfoque documentado prioriza que ELAN cree o inserte anotaciones a partir de JSON. 

La trazabilidad técnica se apoya en logs, estados de job, métricas internas y respuestas estructuradas. En un sistema de inferencia, una falla puede originarse por múltiples causas como un archivo de video inexistente, paquete de modelo inválido, vocabulario incompleto, incompatibilidad de pesos, ausencia de procesamiento CUDA, salida con forma inesperada o error de extracción de keypoints. Sin trazabilidad, estas fallas se vuelven difíciles de diagnosticar. Con trazabilidad, el sistema puede informar un código específico y orientar la corrección sin interrumpir la herramienta principal.

La mantenibilidad depende de que los módulos tengan responsabilidades delimitadas. Separar rutas HTTP, esquemas, servicios, almacenamiento, procesamiento de video y runners permite evolucionar una parte sin reescribir todo el sistema. Este principio es especialmente relevante cuando el proyecto debe incorporar nuevos modelos, cambiar el mecanismo de ejecución o adaptar la comunicación con ELAN. Una arquitectura modular también facilita documentar evidencias, revisar pruebas y justificar decisiones técnicas.

La documentación técnica cumple una función de puente entre implementación y reproducibilidad académica. No basta con que el sistema funcione localmente, también debe poder explicarse, verificarse y mantenerse. Los documentos de fases, contratos de endpoints, ejemplos de solicitudes y descripciones de errores permiten que otro lector comprenda el alcance real del componente. Además, ayudan a diferenciar entre funcionalidades implementadas, funcionalidades planificadas y limitaciones pendientes.

La soberanía de los datos es una consideración adicional en corpus audiovisuales. El procesamiento local evita enviar videos sensibles o material de investigación a servicios externos no controlados. Este aspecto resulta relevante cuando se trabaja con datos de personas, señas, contextos educativos o corpus en construcción. La arquitectura de middleware local favorece que la inferencia se ejecute en el equipo del usuario o en un entorno controlado, aunque la protección efectiva también depende de políticas de almacenamiento, permisos y manejo de archivos que deben documentarse en la metodología y validación.

Por otro lado, la mantenibilidad también exige evitar dependencias implícitas. Si ELAN conoce detalles internos de un modelo, cada cambio de arquitectura obligaría a modificar la herramienta de anotación. Si el modelo asume estructuras internas de ELAN, se reduce la posibilidad de reutilizarlo en otros clientes. El middleware resuelve este problema al establecer una interfaz estable. Mientras el contrato externo se mantenga, los cambios internos pueden limitarse al registro de modelos, runners o postprocesadores.

Finalmente, la evaluación del sistema debe distinguir entre validez técnica y validez lingüística. La validez técnica se relaciona con que el middleware reciba solicitudes, ejecute modelos, controle errores y devuelva respuestas bien formadas. La validez lingüística requiere que los segmentos y glosas sean correctos según criterios especializados. 